\chapter*{先别怕，这份说明书是给零基础准备的}
\addcontentsline{toc}{chapter}{先别怕}

\begin{keyidea}[你不需要先会编程]
本册假设你：知道电脑里有``文件''，会用备份软件，\\
但\textbf{从没}学过哈希、掩码、同调、LFSR……\\
所有怪词会先变成菜市场语言，再配图，再演一遍算法。
\end{keyidea}

\textbf{怎么读}
\begin{enumerate}[nosep]
  \item 先读\textbf{第一部分}（从备份讲起 + 名词图画词典）；
  \item 再看\textbf{第二部分}四种刀法动画片；
  \item 觉得够了就停；想抠细节再读后面的进阶细拆；
  \item 附录有``怪词速查''：遇到单词就翻。
\end{enumerate}

\begin{metaphor}[这本书像什么]
不是考卷，是\textbf{绘本 + 分镜脚本}。\\
红剪刀＝切一刀；绿灯＝条件通过；灰格子＝一个字节。
\end{metaphor}

\vspace{0.5em}
\noindent\textbf{修订}：v1.0 形象版；v1.1 技术点全覆盖；\textbf{v1.2 零基础透讲}。
